% Original pages: 113-120
% Figures: none.
% Uncertainty log: none.

We shall then say that we have an exact sequence of inverse systems. The diagram certainly induces homomorphisms
\[
 0\longrightarrow\varprojlim A_n\longrightarrow\varprojlim B_n\longrightarrow\varprojlim C_n\longrightarrow0
\]
but this sequence is not always exact. However, we have

\noindent\textit{Proposition 10.2.\amtarget{10.2} If $0\longrightarrow\{A_n\}\longrightarrow\{B_n\}\longrightarrow\{C_n\}\longrightarrow0$ is an exact sequence of inverse systems, then the sequence}
\[
 0\longrightarrow\varprojlim A_n\longrightarrow\varprojlim B_n\longrightarrow\varprojlim C_n
\]
\textit{is exact. If $\{A_n\}$ is surjective, then the sequence}
\[
 0\longrightarrow\varprojlim A_n\longrightarrow\varprojlim B_n\longrightarrow\varprojlim C_n\longrightarrow0
\]
\textit{is exact.}

\noindent\textit{Proof.} Let $A=\prod_{n=1}^{\infty}A_n$ and define $d^A:A\longrightarrow A$ by
\[
 d^A((a_n))=(a_n-\theta_{n+1}(a_{n+1})).
\]
Then $\mathop{\rm Ker}d^A\cong\varprojlim A_n$. Define $B,C,d^B,d^C$ similarly. We have a commutative diagram with exact rows
\[
\begin{array}{ccccccccc}
0&\longrightarrow&A&\longrightarrow&B&\longrightarrow&C&\longrightarrow&0\\
&&\big\downarrow d^A&&\big\downarrow d^B&&\big\downarrow d^C&&\\
0&\longrightarrow&A&\longrightarrow&B&\longrightarrow&C&\longrightarrow&0,
\end{array}
\]
and consequently an exact sequence
\[
0\longrightarrow\mathop{\rm Ker}d^A\longrightarrow\mathop{\rm Ker}d^B\longrightarrow\mathop{\rm Ker}d^C
\longrightarrow\mathop{\rm Coker}d^A\longrightarrow\mathop{\rm Coker}d^B\longrightarrow\mathop{\rm Coker}d^C\longrightarrow0.
\]
It remains to prove that if $\{A_n\}$ is surjective, then $d^A$ is surjective. We have to solve the equations
\[
 x_n-\theta_{n+1}(x_{n+1})=a_n\qquad(n\geq1)
\]
for $x_n\in A_n$, where the $a_n\in A_n$ are given. We can do this by induction on $n$, because $\theta_{n+1}$ is surjective. \proofbox

\noindent\textit{Remark.} The group $\mathop{\rm Coker}d^A$ is usually denoted by $\varprojlim{}^1 A_n$: it is the first derived functor of the inverse limit.

\noindent\textit{Corollary 10.3.\amtarget{10.3} Let}
\[
0\longrightarrow G'\longrightarrow G\mathop{\longrightarrow}^{p}G''\longrightarrow0
\]
\textit{be an exact sequence of groups, and let the topology of $G$ be defined by a sequence $(G_n)$. Give $G'$ and $G''$ the topologies defined by the sequences $(G'\cap G_n)$ and $(pG_n)$ respectively. Then the sequence}
\[
0\longrightarrow\widehat{G'}\longrightarrow\widehat G\longrightarrow\widehat{G''}\longrightarrow0
\]
\textit{is exact.}

\noindent\textit{Proof.} Apply \amref{10.2} to the exact sequences
\[
0\longrightarrow G'/(G'\cap G_n)\longrightarrow G/G_n\longrightarrow G''/pG_n\longrightarrow0.
\]
\proofbox

Apply \amref{10.3} with $G'=G_n$, $G''=G/G_n$, the latter being given the discrete topology. Thus $\widehat{G''}=G''$, and we obtain

\noindent\textit{Corollary 10.4.\amtarget{10.4} $\widehat{G_n}$ is a subgroup of $\widehat G$, and $\widehat G/\widehat{G_n}\cong G/G_n$.} \proofbox

Taking inverse limits, we obtain

\noindent\textit{Proposition 10.5.\amtarget{10.5} $\widehat{\widehat G}\cong\widehat G$.} \proofbox

If the canonical mapping $\phi:G\longrightarrow\widehat G$ is an isomorphism, $G$ is said to be \textit{complete}. Thus the completion of any group is complete. Notice that our definition of completeness includes the Hausdorff condition, by \amref{10.1}.

Let $A$ be a ring, $\mathfrak a$ an ideal of $A$. The $\mathfrak a$-adic topology on $A$ is defined by taking $G=A$, $G_n=\mathfrak a^n$. It is Hausdorff if and only if $\bigcap\mathfrak a^n=(0)$. The completion $\widehat A$ is a topological ring, and the canonical mapping $\phi:A\longrightarrow\widehat A$ is a continuous ring homomorphism, whose kernel is $\bigcap\mathfrak a^n$.

More generally, if $M$ is an $A$-module, its $\mathfrak a$-adic topology is defined by taking $G=M$, $G_n=\mathfrak a^nM$. The completion $\widehat M$ is a topological $\widehat A$-module, the multiplication $\widehat A\times\widehat M\longrightarrow\widehat M$ being defined in the obvious way. If $f:M\longrightarrow N$ is an $A$-module homomorphism, it is continuous for the $\mathfrak a$-adic topologies and therefore defines a homomorphism $\widehat f:\widehat M\longrightarrow\widehat N$.

\noindent\textbf{Examples.} 1) Let $A=k[x]$, $\mathfrak a=(x)$. Then $\widehat A=k[[x]]$, the ring of formal power series in $x$ over $k$.

2) Let $A=\mathbb Z$, $\mathfrak a=(p)$, where $p$ is a prime. Then $\widehat A$ is the ring of $p$-adic integers. Its elements can be written uniquely in the form
\[
\sum_{n=0}^{\infty}a_np^n\qquad(0\leq a_n\leq p-1).
\]
In $\widehat A$, $p^n\longrightarrow0$ as $n\longrightarrow\infty$.

\subsection{Filtrations}

Let $A$ be a ring, $M$ an $A$-module, $\mathfrak a$ an ideal of $A$. A decreasing sequence of submodules
\[
M=M_0\supseteq M_1\supseteq\cdots\supseteq M_n\supseteq\cdots
\]
is called a \textit{filtration} of $M$. It is an $\mathfrak a$-filtration if $\mathfrak aM_n\subseteq M_{n+1}$ for all $n$, and a \textit{stable} $\mathfrak a$-filtration if $\mathfrak aM_n=M_{n+1}$ for all sufficiently large $n$. The sequence $(\mathfrak a^nM)$ is a stable $\mathfrak a$-filtration.

\noindent\textit{Lemma 10.6.\amtarget{10.6} Any two stable $\mathfrak a$-filtrations $(M_n)$, $(M'_n)$ of $M$ have bounded difference: that is, there exists an integer $n_0$ such that}
\[
M_{n+n_0}\subseteq M'_n,\qquad M'_{n+n_0}\subseteq M_n
\]
\textit{for all $n\geq0$. Hence the two filtrations define the same topology on $M$.}

\noindent\textit{Proof.} It is enough to prove this when $M'_n=\mathfrak a^nM$. Since $\mathfrak aM_n\subseteq M_{n+1}$, we have $\mathfrak a^nM\subseteq M_n$. Also, for some $n_0$, $M_{n+1}=\mathfrak aM_n$ for all $n\geq n_0$, and hence
\[
M_{n+n_0}=\mathfrak a^nM_{n_0}\subseteq\mathfrak a^nM.
\]
\proofbox

\subsection{Graded rings and modules}

A \textit{graded ring} is a ring $A$, together with a family $(A_n)_{n\geq0}$ of subgroups of its additive group, such that
\[
A=\bigoplus_{n=0}^{\infty}A_n,\qquad A_mA_n\subseteq A_{m+n}.
\]
In particular, $A_0$ is a subring of $A$, and each $A_n$ is an $A_0$-module.

\noindent\textbf{Example.} $A=k[x_1,\ldots,x_r]$, $A_n=$ set of all homogeneous polynomials of degree $n$.

If $A$ is a graded ring, a \textit{graded $A$-module} is an $A$-module $M$, together with a family $(M_n)_{n\geq0}$ of subgroups of its additive group, such that
\[
M=\bigoplus_{n=0}^{\infty}M_n,\qquad A_mM_n\subseteq M_{m+n}.
\]
The elements of $A_n$ (or $M_n$) are called \textit{homogeneous} of degree $n$. Each element $y$ of a graded module $M$ has a unique expression $y=\sum_ny_n$, where $y_n\in M_n$ and all but a finite number of the $y_n$ are zero. The non-zero $y_n$ are called the \textit{homogeneous components} of $y$.

A homomorphism $f:M\longrightarrow N$ of graded $A$-modules is a graded homomorphism if $f(M_n)\subseteq N_n$ for all $n$.

The ideal $A_+=\bigoplus_{n>0}A_n$ is called the \textit{irrelevant ideal} of $A$.

\noindent\textit{Proposition 10.7.\amtarget{10.7} The following conditions are equivalent:}
\begin{enumerate}
\renewcommand{\labelenumi}{\roman{enumi})}
\item \textit{$A$ is Noetherian;}
\item \textit{$A_0$ is Noetherian and $A$ is a finitely-generated $A_0$-algebra.}
\end{enumerate}

\noindent\textit{Proof.} i) $\Rightarrow$ ii): $A_0\cong A/A_+$ is Noetherian, and the ideal $A_+$ is finitely generated. We may take a finite set of homogeneous generators $x_1,\ldots,x_s$, of degrees $k_1,\ldots,k_s>0$. Let $A'$ be the subring of $A$ generated by $A_0$ and the $x_i$. We shall prove by induction on $n$ that $A_n\subseteq A'$. Let $y\in A_n$. Then $y=\sum_i a_ix_i$, where $a_i\in A_{n-k_i}$ (we put $A_m=0$ if $m<0$). By the inductive hypothesis, each $a_i\in A'$, and hence $y\in A'$. Thus $A=A'$.

ii) $\Rightarrow$ i) by Hilbert's basis theorem \amref{7.6}. \proofbox

Let $A$ now be any ring (not necessarily graded), and $\mathfrak a$ an ideal of $A$. Then
\[
A^*=\bigoplus_{n=0}^{\infty}\mathfrak a^n
\]
is a graded ring. If $M$ is an $A$-module, and $(M_n)$ is an $\mathfrak a$-filtration of $M$, then
\[
M^*=\bigoplus_{n=0}^{\infty}M_n
\]
is a graded $A^*$-module, because $\mathfrak a^mM_n\subseteq M_{m+n}$.

If $A$ is Noetherian, and $\mathfrak a$ is generated by $x_1,\ldots,x_r$, then $A^*=A[x_1,\ldots,x_r]$, and is therefore Noetherian.

\noindent\textit{Lemma 10.8.\amtarget{10.8} Let $A$ be Noetherian, $M$ a finitely-generated $A$-module, $(M_n)$ an $\mathfrak a$-filtration of $M$. Then the following are equivalent:}
\begin{enumerate}
\renewcommand{\labelenumi}{\roman{enumi})}
\item \textit{$M^*$ is a finitely-generated $A^*$-module;}
\item \textit{the filtration $(M_n)$ is stable.}
\end{enumerate}

\noindent\textit{Proof.} Each $M_n$ is finitely generated, since $A$ is Noetherian. Let $Q_n=\bigoplus_{r=0}^nM_r$ (a subgroup, but not in general an $A^*$-submodule, of $M^*$). The $A^*$-submodule generated by $Q_n$ is
\[
M_n^*=M_0\oplus\cdots\oplus M_n\oplus\mathfrak aM_n\oplus\mathfrak a^2M_n\oplus\cdots\oplus\mathfrak a^rM_n\oplus\cdots,
\]
and is finitely generated, because $Q_n$ is finitely generated as an $A$-module. The $M_n^*$ form an ascending sequence whose union is $M^*$. Since $A^*$ is Noetherian, $M^*$ is finitely generated if and only if this sequence stops, if and only if $M^*=M_{n_0}^*$ for some $n_0$, if and only if
\[
M_{n_0+r}=\mathfrak a^rM_{n_0}\qquad(r\geq0),
\]
if and only if the filtration is stable. \proofbox

\noindent\textit{Proposition 10.9.\amtarget{10.9} (Artin--Rees lemma). Let $A$ be a Noetherian ring, $\mathfrak a$ an ideal of $A$, $M$ a finitely-generated $A$-module, $(M_n)$ a stable $\mathfrak a$-filtration of $M$. Let $M'$ be a submodule of $M$. Then $(M'\cap M_n)$ is a stable $\mathfrak a$-filtration of $M'$.}

\noindent\textit{Proof.} We have
\[
\mathfrak a(M'\cap M_n)\subseteq\mathfrak aM'\cap\mathfrak aM_n\subseteq M'\cap M_{n+1},
\]
so that $(M'\cap M_n)$ is an $\mathfrak a$-filtration. The corresponding graded module is a graded submodule of $M^*$, and is therefore finitely generated, since $A^*$ is Noetherian. Hence the filtration is stable, by \amref{10.8}. \proofbox

Taking $M_n=\mathfrak a^nM$, we obtain

\noindent\textit{Corollary 10.10.\amtarget{10.10} There exists an integer $k$ such that}
\[
(\mathfrak a^nM)\cap M'=\mathfrak a^{n-k}((\mathfrak a^kM)\cap M')\qquad(n\geq k).
\]
\proofbox

From \amref{10.6} we therefore obtain

\noindent\textit{Theorem 10.11.\amtarget{10.11} Let $A$ be a Noetherian ring, $\mathfrak a$ an ideal, $M$ a finitely-generated $A$-module, $M'$ a submodule. Then the filtrations $(\mathfrak a^nM')$ and $(\mathfrak a^nM\cap M')$ have bounded difference. In particular, the $\mathfrak a$-topology of $M'$ is the same as the topology induced by the $\mathfrak a$-topology of $M$.} \proofbox

\noindent\textit{Remark.} In this chapter we shall use only the topological part of \amref{10.11}. In the next chapter we shall need the stronger statement that the filtrations have bounded difference.

\noindent\textit{Proposition 10.12.\amtarget{10.12} Let}
\[
0\longrightarrow M'\longrightarrow M\longrightarrow M''\longrightarrow0
\]
\textit{be an exact sequence of finitely-generated modules over a Noetherian ring $A$, and let $\mathfrak a$ be an ideal of $A$. Then the sequence}
\[
0\longrightarrow\widehat{M'}\longrightarrow\widehat M\longrightarrow\widehat{M''}\longrightarrow0
\]
\textit{is exact.} \proofbox

There is a natural homomorphism $A\longrightarrow\widehat A$, which makes $\widehat A$ an $A$-algebra. We shall now compare $\widehat A\otimes_AM$ and $\widehat M$. There is a natural homomorphism
\[
\widehat A\otimes_AM\longrightarrow\widehat A\otimes_A\widehat M\longrightarrow\widehat A\otimes_{\widehat A}\widehat M=\widehat M.
\]

\noindent\textit{Proposition 10.13.\amtarget{10.13} Let $A$ be any ring, $\mathfrak a$ any ideal of $A$, $M$ a finitely-generated $A$-module. Then the natural homomorphism}
\[
\widehat A\otimes_AM\longrightarrow\widehat M
\]
\textit{is surjective. If $A$ is Noetherian, it is an isomorphism.}

\noindent\textit{Proof.} Completion commutes with finite direct sums. Hence, if $F\cong A^n$, then $\widehat A\otimes_AF\cong\widehat F$. Choose an exact sequence
\[
0\longrightarrow N\longrightarrow F\longrightarrow M\longrightarrow0
\]
with $F$ a finitely-generated free module. We have a commutative diagram
\[
\begin{array}{ccccccccc}
\widehat A\otimes_AN&\longrightarrow&\widehat A\otimes_AF&\longrightarrow&\widehat A\otimes_AM&\longrightarrow&0\\
\big\downarrow\gamma&&\big\downarrow\beta&&\big\downarrow\alpha&&\\
\widehat N&\longrightarrow&\widehat F&\mathop{\longrightarrow}^{\delta}&\widehat M&\longrightarrow&0.
\end{array}
\]
The top row is exact by \amref{2.18}. By \amref{10.3}, $\delta$ is surjective. Since $\beta$ is an isomorphism, it follows that $\alpha$ is surjective. If $A$ is Noetherian, then $N$ is finitely generated, so that $\gamma$ is surjective, and the bottom row is exact by \amref{10.12}. A diagram chase now shows that $\alpha$ is injective. \proofbox

Propositions \amref{10.12} and \amref{10.13} show that, if $A$ is Noetherian, the functor $M\mapsto\widehat A\otimes_AM$ is exact on the category of finitely-generated $A$-modules. Hence

\noindent\textit{Proposition 10.14.\amtarget{10.14} If $A$ is Noetherian, $\widehat A$ is a flat $A$-algebra.} \proofbox

\noindent\textit{Remark.} If $M$ is not finitely generated, the functor $M\mapsto\widehat M$ is not in general exact. The ``good'' functor is $M\mapsto\widehat A\otimes_AM$, which is exact for all modules $M$ and coincides with $M\mapsto\widehat M$ when $M$ is finitely generated.

\noindent\textit{Proposition 10.15.\amtarget{10.15} Let $A$ be Noetherian, $\mathfrak a$ an ideal of $A$. Then}
\begin{enumerate}
\renewcommand{\labelenumi}{\roman{enumi})}
\item \textit{$\widehat{\mathfrak a}=\widehat A\mathfrak a\cong\widehat A\otimes_A\mathfrak a$;}
\item \textit{$\widehat{\mathfrak a^n}=(\widehat{\mathfrak a})^n$;}
\item \textit{$\mathfrak a^n/\mathfrak a^{n+1}\cong\widehat{\mathfrak a}^{n}/\widehat{\mathfrak a}^{n+1}$;}
\item \textit{$\widehat{\mathfrak a}$ is contained in the Jacobson radical of $\widehat A$.}
\end{enumerate}

\noindent\textit{Proof.} Since $\mathfrak a$ is finitely generated, the natural homomorphism $\widehat A\otimes_A\mathfrak a\longrightarrow\widehat{\mathfrak a}$ is an isomorphism, by \amref{10.13}. Its image is $\widehat A\mathfrak a$, which proves i). Applying i) to $\mathfrak a^n$, we have
\[
\widehat{\mathfrak a^n}=\widehat A\mathfrak a^n=(\widehat A\mathfrak a)^n
\]
by \amref{1.18}, and this is equal to $(\widehat{\mathfrak a})^n$ by i). This proves ii). By \amref{10.4},
\[
A/\mathfrak a^n\cong\widehat A/\widehat{\mathfrak a}^{n},
\]
and iii) follows on taking quotients. By ii) and \amref{10.5}, $\widehat A$ is complete for the $\widehat{\mathfrak a}$-adic topology. Hence, if $x\in\widehat{\mathfrak a}$, the series $1+x+x^2+\cdots$ converges in $\widehat A$, and its sum is the inverse of $1-x$. Thus $\widehat{\mathfrak a}$ is contained in the Jacobson radical of $\widehat A$. \proofbox

\noindent\textit{Proposition 10.16.\amtarget{10.16} Let $A$ be a Noetherian local ring, $\mathfrak m$ its maximal ideal. Then the $\mathfrak m$-adic completion $\widehat A$ of $A$ is a local ring, its maximal ideal being $\widehat{\mathfrak m}$.}

\noindent\textit{Proof.} $\widehat A/\widehat{\mathfrak m}\cong A/\mathfrak m$ is a field, and $\widehat{\mathfrak m}$ is contained in the Jacobson radical of $\widehat A$: hence $\widehat{\mathfrak m}$ is the unique maximal ideal of $\widehat A$. \proofbox

\subsection{Krull's theorem}

We can now determine the kernel of the canonical homomorphism $M\longrightarrow\widehat M$.

\noindent\textit{Theorem 10.17.\amtarget{10.17} Let $A$ be a Noetherian ring, $\mathfrak a$ an ideal of $A$, $M$ a finitely-generated $A$-module, and let $\widehat M$ be the $\mathfrak a$-adic completion of $M$. Then the kernel $E=\bigcap_{n=1}^{\infty}\mathfrak a^nM$ of the canonical mapping $M\longrightarrow\widehat M$ consists of those elements of $M$ which are annihilated by some element of $1+\mathfrak a$.}

\noindent\textit{Proof.} The topology induced on $E$ by the $\mathfrak a$-topology of $M$ is the trivial topology; hence by \amref{10.11} the $\mathfrak a$-topology of $E$ is the trivial topology. Thus $\mathfrak aE=E$. Since $E$ is finitely generated, it follows from \amref{2.5} that $(1-a)E=0$ for some $a\in\mathfrak a$. Conversely, if $(1-a)x=0$, where $a\in\mathfrak a$, then
\[
x=ax=a^2x=\cdots\in\bigcap_{n=1}^{\infty}\mathfrak a^nM=E.
\]
\proofbox

\noindent\textit{Remarks.} 1) Let $S=1+\mathfrak a$. The canonical mappings $A\longrightarrow\widehat A$ and $A\longrightarrow S^{-1}A$ have the same kernel. For each $a\in\mathfrak a$, the series
\[
(1-a)^{-1}=1+a+a^2+\cdots
\]
converges in $\widehat A$. Hence the elements of $S$ are units in $\widehat A$, and there is a natural injective homomorphism $S^{-1}A\longrightarrow\widehat A$.

2) The theorem fails if $A$ is not Noetherian. For example, let $A$ be the ring of $C^{\infty}$ functions on the real line, and let $\mathfrak a$ be the ideal of functions which vanish at the origin. Then $\mathfrak a$ is a maximal ideal, since $A/\mathfrak a\cong\mathbb R$. The ideal $\mathfrak a$ is generated by the identity function $x$, and $\bigcap\mathfrak a^n$ is the ideal of all functions which are \textit{flat} at the origin (i.e., all their derivatives vanish at the origin). A function $f$ is annihilated by some element of $1+\mathfrak a$ if and only if $f$ vanishes identically in some neighborhood of the origin. Thus the function
\[
f(x)=\begin{cases}e^{-1/x^2},&x\neq0,\\0,&x=0,\end{cases}
\]
is a counterexample. In this example the kernels of $A\longrightarrow\widehat A$ and $A\longrightarrow S^{-1}A$ $(S=1+\mathfrak a)$ are different. The ring $A$ is not Noetherian.

\noindent\textit{Corollary 10.18.\amtarget{10.18} Let $A$ be a Noetherian integral domain, $\mathfrak a\neq(1)$ an ideal of $A$. Then $\bigcap_{n=1}^{\infty}\mathfrak a^n=0$.}

\noindent\textit{Proof.} The elements of $1+\mathfrak a$ are not zero-divisors. \proofbox

\noindent\textit{Corollary 10.19.\amtarget{10.19} Let $A$ be a Noetherian ring, $\mathfrak a$ an ideal of $A$ contained in the Jacobson radical of $A$; let $M$ be a finitely-generated $A$-module. Then the $\mathfrak a$-topology of $M$ is Hausdorff. In other words, $\bigcap_{n=1}^{\infty}\mathfrak a^nM=0$.}

\noindent\textit{Proof.} By \amref{1.9}, every element of $1+\mathfrak a$ is a unit. \proofbox

\noindent\textit{Corollary 10.20.\amtarget{10.20} Let $A$ be a Noetherian local ring, $M$ a finitely-generated $A$-module. Then the $\mathfrak m$-topology of $M$ is Hausdorff. In particular, the $\mathfrak m$-topology of $A$ is Hausdorff.} \proofbox

We can restate this last result in terms of primary ideals. An $\mathfrak m$-primary ideal is any ideal contained between $\mathfrak m$ and some power $\mathfrak m^n$; hence the intersection of all $\mathfrak m$-primary ideals is zero. More generally, let $A$ be any Noetherian ring, $\mathfrak p$ any prime ideal. Apply \amref{10.20} to the local ring $A_{\mathfrak p}$. There is a one-to-one correspondence between the $\mathfrak p$-primary ideals of $A$ and the $\mathfrak m$-primary ideals of $A_{\mathfrak p}$, where $\mathfrak m=\mathfrak pA_{\mathfrak p}$. Hence

\noindent\textit{Corollary 10.21.\amtarget{10.21} The intersection of all $\mathfrak p$-primary ideals of $A$ is the kernel of the canonical mapping $A\longrightarrow A_{\mathfrak p}$.} \proofbox

\subsection{The associated graded ring}

Let $A$ be a ring, $\mathfrak a$ an ideal of $A$. We define
\[
G(A)\;(=G_{\mathfrak a}(A))=\bigoplus_{n=0}^{\infty}\mathfrak a^n/\mathfrak a^{n+1}\qquad(\mathfrak a^0=A).
\]
We make $G(A)$ into a graded ring as follows. If $x_n\in\mathfrak a^n$, let $\bar x_n$ denote its image in $\mathfrak a^n/\mathfrak a^{n+1}$. We define
\[
\bar x_m\bar x_n=\overline{x_mx_n}\in\mathfrak a^{m+n}/\mathfrak a^{m+n+1}.
\]

More generally, let $M$ be an $A$-module and $(M_n)$ an $\mathfrak a$-filtration of $M$. We define
\[
G(M)=\bigoplus_{n=0}^{\infty}M_n/M_{n+1}.
\]
Then $G(M)$ is a graded $G(A)$-module. We write $G_n(M)=M_n/M_{n+1}$.

\noindent\textit{Proposition 10.22.\amtarget{10.22}}
\begin{enumerate}
\renewcommand{\labelenumi}{\roman{enumi})}
\item \textit{$G_{\mathfrak a}(A)$ is a Noetherian ring.}
\item \textit{$G_{\mathfrak a}(A)$ and $G_{\widehat{\mathfrak a}}(\widehat A)$ are isomorphic graded rings.}
\item \textit{If $M$ is a finitely-generated $A$-module and $(M_n)$ is a stable $\mathfrak a$-filtration of $M$, then $G(M)$ is a finitely-generated graded $G_{\mathfrak a}(A)$-module.}
\end{enumerate}

\noindent\textit{Proof.} i) Let $\mathfrak a$ be generated by $x_1,\ldots,x_s$, and let $\bar x_i$ denote the image of $x_i$ in $\mathfrak a/\mathfrak a^2$. Then
\[
G(A)=(A/\mathfrak a)[\bar x_1,\ldots,\bar x_s],
\]
and the result follows from Hilbert's basis theorem.

ii) By \amref{10.15}, iii), we have
\[
\mathfrak a^n/\mathfrak a^{n+1}\cong\widehat{\mathfrak a}^{n}/\widehat{\mathfrak a}^{n+1}.
\]
